% !TeX program = xelatex
\documentclass[12pt,a4paper]{article}
\usepackage{amsmath,amssymb}
\usepackage{geometry}
\usepackage{booktabs}
\usepackage{graphicx}
\usepackage{xcolor}
\usepackage[UTF8]{ctex}
\geometry{margin=2.5cm}
\title{Q3: 回收体系稳定性分析}
\author{}
\date{}

\begin{document}
\maketitle

\section{模型（来自 Q2）}

对称拮抗模型。$\alpha$（促进）、$\beta$（抑制）\textbf{恒为正数}：
\begin{equation}\label{eq:ode}
\frac{dx}{dt} = \frac{\gamma}{1+\beta} \cdot x \cdot \left(1 - \frac{x}{K(1+\alpha)}\right), \qquad \alpha \geq 0,\; \beta \geq 0
\end{equation}

标准 logistic 形式：
\begin{equation}\label{eq:logistic}
\frac{dx}{dt} = r \cdot x \cdot \left(1 - \frac{x}{K_{\text{eff}}}\right), \qquad
r = \frac{\gamma}{1+\beta}, \quad K_{\text{eff}} = K(1+\alpha)
\end{equation}

拟合结果：$r = 0.008485$, $K_{\text{eff}} = 8142.22$, $\alpha = 0.0164$, $\beta = 1.5929$。

\section{平衡点与稳定性}

令 $dx/dt = 0$：
\begin{equation}
r \cdot x \cdot \left(1 - \frac{x}{K_{\text{eff}}}\right) = 0
\quad\Rightarrow\quad x_1^* = 0,\;\; x_2^* = K_{\text{eff}}
\end{equation}

Jacobi 导数：
\begin{equation}
f'(x) = r\left(1 - \frac{2x}{K_{\text{eff}}}\right)
\end{equation}

\begin{center}
\begin{tabular}{lccc}
\toprule
平衡点 & 值 & 特征值 $\lambda$ & 稳定性 \\
\midrule
$x_1^*$ & $0$ & $\lambda_1 = r$ & \textbf{不稳定}（$r > 0$ 恒成立）\\
$x_2^*$ & $K_{\text{eff}} = K(1+\alpha)$ & $\lambda_2 = -r$ & \textbf{渐近稳定}（$-r < 0$ 恒成立）\\
\bottomrule
\end{tabular}
\end{center}

\section{结构稳定性}

\begin{equation}
\beta \geq 0 \;\Rightarrow\; 1+\beta \geq 1 \;\Rightarrow\; r = \frac{\gamma}{1+\beta} \in (0,\, \gamma]
\end{equation}

\textbf{关键结论}：$\alpha, \beta \geq 0$ 是定义本身（促进系数和抑制系数），
因此在整个有效参数空间中 $r > 0$ \textbf{恒成立}。

\begin{center}
\fbox{%
\begin{minipage}{0.85\textwidth}
\centering
\textbf{系统具有结构稳定性。}\\[3pt]
对任意 $\alpha \geq 0$, $\beta \geq 0$ 和任意正初值 $x(0) > 0$，
恒有 $x(t) \to K_{\text{eff}} > 0$。\\
\textbf{有效参数空间内不存在任何分岔点。}
\end{minipage}}
\end{center}

这意味着``沪尚回收''体系在数学上不存在分岔失稳——只要促进效应
和抑制效应保持其设计定义（正数），回收量将始终趋向正平衡值 $K_{\text{eff}}$。

\section{实用安全边界}

虽然系统在数学上全局稳定，但运行效率要求参数保持在合理范围内。
安全边界由\textbf{性能约束}定义，而非分岔：

\subsection{增长率约束}

\begin{equation}
r = \frac{\gamma}{1+\beta} \geq r_{\min} \quad\Rightarrow\quad
\beta \leq \frac{\gamma}{r_{\min}} - 1
\end{equation}

$r_{\min}$ 的取值依据为：Logistic 系统恢复至 $99\%$ 饱和水平所需时间满足保守上界
$t_{99\%} \leq -\ln(0.01)/r = \ln(100)/r$。要求 $t_{99\%}$ 不超过
$2.5$ 个观测周期（$2.5 \times 365 = 912.5$ 天），即
$r \geq \ln(100)/912.5 \approx 0.00505$，取整得 $r_{\min} = 0.005$。
代入得 $\beta_{\max} = 0.022/0.005 - 1 = 3.40$。

\subsection{容量约束}

\begin{equation}
K_{\text{eff}} = K(1+\alpha) \geq K_{\min} \quad\Rightarrow\quad
\alpha \geq \frac{K_{\min}}{K} - 1
\end{equation}

取 $K_{\min} = K$：$\alpha \geq 0$——由定义自然满足。

\subsection{安全运行范围}

\begin{center}
\begin{tabular}{lccc}
\toprule
参数 & 下界 & 上界 & 当前值 \\
\midrule
$\alpha$（促进） & $0$（定义） & — & $0.0164$（容量为 $K$ 的 $101.6\%$）\\
$\beta$（抑制） & $0$（定义） & $3.40$（$r_{\min}$ 约束） & $1.59$（$r$ 为 $\gamma$ 的 $38.6\%$）\\
\bottomrule
\end{tabular}
\end{center}

\begin{center}
\begin{tabular}{lc}
\toprule
安全裕度 & 数值 \\
\midrule
$\beta$ 距 $\beta_{\max}$ & $1.81$ \\
$r$ 高于 $r_{\min}$ & $0.0035$ \\
$r/\gamma$ & $38.6\%$ \\
\bottomrule
\end{tabular}
\end{center}

\section{解耦调控}

该模型的一个重要优势：$\alpha$ 和 $\beta$ \textbf{完全解耦}。

\begin{itemize}
  \item $\alpha$ \textbf{仅影响} $K_{\text{eff}}$：调节促进力度可改变饱和容量而不影响收敛速度
  \item $\beta$ \textbf{仅影响} $r$：控制抑制力度可改变收敛速度而不影响饱和容量
\end{itemize}

这一特性使得决策者能够\textbf{独立调控}两个维度：通过扩大站点覆盖提升总容量（$\alpha$），
同时通过改善运维效率维持增长速度（降低 $\beta$）。

\section{关键结论}

\begin{enumerate}
  \item \textbf{结构稳定}：$\alpha, \beta \geq 0$ 是定义，$r = \gamma/(1+\beta) > 0$ 恒成立。
        整个有效参数空间内无分岔点。
  \item \textbf{设计保证安全}：只要促进效应和抑制效应不违背其定义（恒为正），
        系统不会发生分岔失稳。
  \item \textbf{风险在于性能退化}：$\beta$ 过高导致 $r$ 过小（收敛过慢）是实际运营风险，
        而非数学分岔。
  \item \textbf{解耦控制}：$\alpha \to K_{\text{eff}}$（容量）、$\beta \to r$（速度），
        互不干扰，可独立优化。
  \item \textbf{安全边界}：$\beta \leq \gamma/r_{\min} - 1 = 3.40$（$r_{\min} = 0.005$），
        当前 $\beta = 1.59$ 裕度充足。
\end{enumerate}

\end{document}
